\newpage

\section{Введение}

Алгоритм Frank--Wolfe, также известный как метод условного градиента, был предложен для задач квадратичного программирования и стал одним из базовых проекционно-свободных методов выпуклой оптимизации~\cite{frank1956algorithm}.
В отличие от проекционных градиентных методов, на каждой итерации он решает линейную задачу минимизации на допустимом множестве.
Это делает метод особенно удобным для задач, в которых линейный оракул существенно дешевле евклидовой проекции: оптимизация на симплексах, транспортных политопах, ядерных шарах, множествах с ограничением на $l_1$-норму и других структурированных областях~\cite{jaggi2013revisiting,braun2022conditional}.

В машинном обучении оптимизация параметров обычно формулируется как минимизация функции потерь с регуляризацией.
Регуляризованная задача
\[
    f(x) + \mu g(x) \to \min_x
\]
может быть рассмотрена как задача с ограничением
\[
    f(x) \to \min_{x \in \cX},
\]
где множество $\cX$ задает допустимую область параметров.
Если для $\cX$ эффективен линейный оракул, методы Frank--Wolfe позволяют решать такие задачи без явного вычисления проекции.
Это особенно важно в задачах высокой размерности, где проекция может быть основной вычислительной нагрузкой.

Классический алгоритм Frank--Wolfe эффективен на первых итерациях, но часто замедляется при приближении к решению.
Одной из причин является зигзагообразное движение траектории около условного минимума~\cite{patriksson2015traffic}.
Однако зигзагообразность является только одним из вычислительных ограничений метода.
В разных постановках узким местом может быть не только геометрия направления, но и стоимость
полного линейного оракула, вычисление информации второго порядка или коммуникация между
агентами.
Поэтому развитие Frank--Wolfe естественно идет не по одной линии, а по нескольким
взаимосвязанным исследовательским направлениям.

Первое направление удобно анализировать в бекмановской постановке равновесного распределения
транспортных потоков.
В этой задаче линейный оракул имеет понятную структуру кратчайших путей, а гессиан целевой
функции диагонален, что позволяет детально исследовать модификации направления шага:
CFW, BFW, NFW, FFW и WFFW.
Второе направление переносит идею сопряженных направлений в задачи машинного обучения.
Здесь явное использование гессиана и точных Hessian-vector products становится
вычислительно дорогим, поэтому произведения гессиана на направления аппроксимируются
конечными разностями градиентов.
Третье направление снова возникает в транспортных сетях, но отвечает на другой вопрос:
как уменьшить стоимость полного LMO, если на каждой итерации приходится решать задачи
кратчайших путей для большого числа источников.
Это приводит к SOFW как стохастическому блочно-координатному варианту FW.
Наконец, четвертое направление рассматривает распределенную постановку, где у агентов есть
локальные функции, а алгоритм должен совмещать блочный линейный оракул, консенсус и
отслеживание градиента.
В этой части DBCFW выступает как распределенное обобщение блочно-координатного подхода:
блочная структура сохраняет дешевизну локальных LMO-шагов, а механизмы консенсуса и
отслеживания градиента позволяют согласовывать локальные обновления между агентами.

Таким образом, работа объединяет четыре связанных направления модификации Frank--Wolfe:
\begin{enumerate}
    \item модификации FW в задаче равновесного распределения транспортных потоков;
    \item перенос Conjugate Frank--Wolfe и N-conjugate Frank--Wolfe на задачи машинного обучения;
    \item стохастический блочно-координатный Frank--Wolfe для крупномасштабных транспортных сетей;
    \item децентрализованный блочно-координатный Frank--Wolfe для распределенной оптимизации.
\end{enumerate}

\subsection{Актуальность}

Актуальность работы связана с тем, что многие современные задачи машинного обучения
и сетевой оптимизации имеют одновременно три свойства.
Во-первых, размерность переменной велика, и поэтому проекция на допустимое множество
может быть дорогой операцией.
Во-вторых, допустимое множество часто имеет структуру, для которой линейный оракул
вычисляется существенно проще проекции.
В-третьих, данные и вычисления могут быть распределены: либо по блокам задачи, либо по
агентам, либо по источникам транспортного спроса.

Алгоритм Frank--Wolfe естественно подходит для таких условий, потому что он заменяет
проекцию линейной минимизацией.
Однако базовый FW не решает все вычислительные проблемы.
В транспортных задачах он может требовать слишком дорогой полной итерации.
В задачах с активной границей допустимого множества он часто демонстрирует зигзагообразное
движение.
В распределенных задачах необходимо дополнительно учитывать консенсус между агентами и
ошибку приближения глобального градиента.
Поэтому практическая эффективность FW зависит не только от классической теоретической
оценки сходимости, но и от архитектуры итерации.

Рассматриваемые в работе методы образуют единую исследовательскую линию.
Первое направление изучает геометрию шага FW и использует память о предыдущих направлениях
для подавления зигзагообразного поведения.
Второе направление адаптирует эту идею к задачам машинного обучения, где информация второго
порядка должна извлекаться без явного построения гессиана.
Третье направление рассматривает стоимость полного линейного оракула как самостоятельное
вычислительное ограничение и приводит к стохастическим блочно-координатным схемам.
Четвертое направление объединяет блочный оракул с децентрализованной оптимизацией,
консенсусом и отслеживанием градиента.

\subsection{Научная новизна и практическая значимость}

Научная новизна работы состоит в развитии нескольких модификаций Frank--Wolfe как единой
линии проекционно-свободных методов с явным теоретическим обоснованием.
CFW, NFW, SOFW и DBCFW обычно можно рассматривать отдельно: как транспортные методы,
как методы машинного обучения, как стохастические блочные схемы или как распределенную
оптимизацию.
В этой работе они рассматриваются как разные способы устранения вычислительных ограничений
одного подхода: неблагоприятной геометрии направления, дорогой информации второго порядка,
дорогого полного линейного оракула и ограниченной коммуникации между агентами.

Теоретическая часть работы включает три ключевых результата.
Во-первых, для CFW в общей гладкой выпуклой постановке доказывается оценка сходимости по
минимальному зазору Frank--Wolfe.
Эта оценка показывает, что переход к сопряженным направлениям сохраняет базовый порядок
сходимости FW, а отличие входит только через контролируемый параметр направления.
Во-вторых, для SOFW формулируется стохастический блочно-координатный LMO-оракул и
доказывается оценка ожидаемой субоптимальности, в том числе для батчевого выбора блоков.
Тем самым показывается, что можно уменьшать стоимость итерации за счет случайного
обновления части блоков без ухудшения базового порядка теоретической гарантии.
В-третьих, для DBCFW выводится рекурсия спуска и оценка сходимости для распределенной
блочно-координатной схемы с консенсусом и отслеживанием градиента.
Эта оценка показывает, при каких условиях децентрализованный блочный вариант сохраняет
порядок $O(1/t)$ для гладких выпуклых задач.

Практическая значимость состоит в том, что такая архитектура помогает выбирать модификацию
под вычислительный узкий участок задачи.
Если проблема заключается в зигзагообразном движении при активном ограничении, естественно
использовать CFW или NFW.
Если основной расход связан с полным линейным оракулом по большому числу блоков, подходит
BCFW/SOFW.
Если данные распределены между агентами и при этом множество имеет блочную структуру,
естественным кандидатом становится DBCFW.

Кроме того, перенос CFW и NFW в машинное обучение показывает важный общий принцип.
Методы можно адаптировать к общей гладкой задаче без явного построения гессиана:
произведения гессиана на направления приближаются конечными разностями градиентов.
Это сохраняет идею сопряженности, но оставляет итерацию достаточно дешевой для практического
обучения моделей.

\subsection{Цель работы}

Целью работы является систематизация, развитие и теоретическое обоснование модификаций
алгоритма Frank--Wolfe для задач машинного обучения и смежных крупномасштабных задач
выпуклой оптимизации, в которых вычислительная стоимость линейного оракула, выбора
направления и коммуникаций существенно влияет на практическую эффективность метода.
Особое внимание уделяется получению оценок сходимости для CFW, SOFW и DBCFW, поскольку
именно эти результаты показывают, что ускорение отдельных компонентов итерации не разрушает
базовые гарантии метода.

\subsection{Задачи работы}

\begin{enumerate}
    \item Описать базовую постановку задачи равновесного распределения транспортных потоков и модификации FW, использующие предыдущие направления.
    \item Сформулировать CFW и NFW для общей гладкой выпуклой задачи на компактном множестве и показать, как коэффициенты сопряженности можно вычислять без явного построения гессиана.
    \item Доказать теоретическую оценку сходимости CFW по зазору Frank--Wolfe и показать, что порядок сходимости сохраняется относительно базового FW.
    \item Описать SOFW как стохастический блочно-координатный вариант FW, обновляющий только часть блоков корреспонденций, и показать его связь с Block-Coordinate Frank--Wolfe.
    \item Получить оценку ожидаемой сходимости SOFW, включая батчевый выбор блоков, и проанализировать компромисс между стоимостью итерации, дисперсией направления и качеством теоретической гарантии.
    \item Сформулировать Decentralized Block-Coordinate Frank--Wolfe как распределенное обобщение блочно-координатных FW-методов с консенсусом и отслеживанием градиента.
    \item Доказать рекурсию спуска и оценку сходимости DBCFW, показать, что DBCFW содержит BCFW и децентрализованный FW как специальные случаи, и описать SOFW в том же блочно-координатном языке.
\end{enumerate}

\subsection{Структура работы}

Глава~\ref{sec:fw-modifications} посвящена модификациям Frank--Wolfe в модели Бекмана.
Глава~\ref{sec:cfw-ml} переносит CFW и NFW в постановку машинного обучения и содержит
доказательство оценки сходимости CFW.
Глава~\ref{sec:scfw} рассматривает стохастический блочно-координатный подход для крупных
транспортных графов и выводит оценку ожидаемой сходимости SOFW.
Глава~\ref{sec:dbcfw} формулирует децентрализованный блочно-координатный алгоритм, его
связь с предыдущими методами и теоретическую оценку сходимости.
